Una de les llunes de Júpiter, Io, descriu una òrbita de radi mitjà $4,22\times 10^{8}\ \mathrm{m}$ i de període $1,53\times 10^{5}\ \mathrm{s}$.

\begin{apartat}{1,25}
Calculeu la massa de Júpiter.
\end{apartat}

\begin{apartat}{1,25}
Calculeu el radi mitjà de l'òrbita d'una altra lluna de Júpiter, Cal·listo, que té un període d'$1,44\times 10^{6}\ \mathrm{s}$.
\end{apartat}

\blocetiquetat{Dada:}{$G = 6,67\times 10^{-11}\ \mathrm{N\,m^2\,kg^{-2}}$.}
